%%%%%%%%%%  scan idx 22  |  folio 9  %%%%%%%%%%
Si $f$ est un
morphisme localement d'intersection compl\`ete, on introduit une notion naturelle de ``dimension relative virtuelle'' $d(f)$ par la formule
\[
d(f)(x)=\dim_x\bigl(f'^{-1}f(x)\bigr)-\codim_x(X,X').
\]
C'est une fonction \`a valeurs enti\`eres localement constante en le point $x\in X$, donc constante si $X$ est connexe, qui ne d\'epend d'ailleurs pas non plus du choix de la factorisation (2.1). Lorsque cette dimension relative virtuelle $d$ est constante, que $f$ est propre et que $X$ et $Y$ admettent des Modules inversibles amples, on peut montrer (Exp. VIII), encore que ce soit loin d'\^etre un r\'esultat trivial, que l'on a
\begin{equation}\tag{3.4}
f_*\bigl(\operatorname{Filt}^iK^\bullet(X)_{\mathbb Q}\bigr)\subset \operatorname{Filt}^{i-d}K^\bullet(Y)_{\mathbb Q},
\end{equation}
d'o\`u un homomorphisme (3.3) de degr\'e $-d$.

\subsubsection*{3.3.}
Nous avons donc r\'euni tous les \'el\'ements de structure pour donner un sens \`a la formule (1.2) \uline{dans le cas d'un morphisme propre et localement d'intersection compl\`ete de sch\'emas noeth\'eriens $X,Y$ admettant tous deux des Modules inversibles amples}. Moyennant un peu de soin suppl\'ementaire, cette formulation garde d'ailleurs un sens ind\'ependamment de toute hypoth\`ese noeth\'erienne.

Sous cette forme g\'en\'erale, cette formule de Riemann--Roch sera prouv\'ee dans Exp. VIII. Le lecteur constatera que la d\'emonstration donn\'ee dans ce S\'eminaire combine des \'el\'ements des deux d\'emonstrations qui figurent dans [RRR] et dont il a \'et\'e question dans l'Introduction. La n\'ecessit\'e de recourir aux m\'ethodes de la premi\`ere de ces d\'emonstrations, impliquant plus de $\lambda$-formalisme que la deuxi\`eme (suivie dans Borel--Serre), est due \`a la n\'ecessit\'e d'une d\'emonstration de l'inclusion (3.4), question qui ne se posait pas dans le cadre plus restreint de [RRR] et Borel--Serre [2], o\`u on pouvait travailler avec l'anneau de Chow ou l'anneau $\Gr^\bullet_{\mathrm{top}}$. De la deuxi\`eme d\'emonstration, nous conservons l'introduction du sch\'ema \'eclat\'e de $Y$ le long de $X$ dans le cas o\`u $f$ est une immersion. Cet artifice technique dispara\^itra sans doute, et la d\'emonstration du cas d'une immersion s'\'etendra au cas o\`u on ne postule plus l'existence d'un Module inversible ample sur $Y$, une fois r\'esolue la question soulev\'ee dans Exp.\ XIV n\textsuperscript{o}~1 sur l'introduction d'op\'erations de puissances ext\'erieures dans la cat\'egorie d\'eriv\'ee $D^-(Y)$.

\subsection*{4.}
Il nous faut enfin donner quelques indications sur la fa\c{c}on dont on peut encore donner un sens \`a (1.2) en l'absence d'une hypoth\`ese d'existence 
